   
\chapter[BSS users and their mobility patterns]{\ac{bss} users and their mobility patterns}
\chaptermark{BSS users and their mobility patterns}
\label{ch:bss_users_and_their_mobility_patterns}

\acp{bss} only work when people actually use them, and how they do so varies widely across cities and user groups. To understand the role of these systems in everyday mobility, it is useful to look at who the users are, how they ride, and what shapes their decisions. This Chapter introduces these aspects by outlining the main user profiles, the factors that influence \ac{bss} usage, and the typical patterns that emerge in where and when people ride shared bicycles.

\textcolor{blue}{
    This Chapter is therefore preparatory for the contributions of this thesis. Part~\ref{pt:data-driven-station-planning} considers where stations should be located and who is able to reach them, which requires an understanding that demand and access are not evenly distributed across space. Part~\ref{pt:predictive_operations} explores how bicycles are used, how batteries drain, and how components experience wear, all of which depend on the heterogeneity of trips, users, and the mix of mechanical and electric bikes discussed below.
}

\section[Who uses BSSs?]{Who uses \acp{bss}?}
\label{sec:who_uses_bss}

Since the early empirical studies of the 2010s, \ac{bss} users have consistently been portrayed as having a fairly narrow socio-demographic profile, typically described as younger, white, male, well-educated, and more affluent than the general population. Early analyses revealed these patterns. In Dublin, 78\% of users were men, 59\% were aged between 25 and 36 years, and a majority reported middle- to high-income levels~\cite{Murphy2015}. In Montreal, 69\% of users were men, with an average age of 35.7 years~\cite{BachandMarleau2012}. In London, over 80\% of trips were made by men, only about 10\% of users came from highly deprived areas~\cite{Goodman2014}, and 78\% of total travel time was undertaken by users under 45 years old~\cite{Woodcock2014}. Ethnic disparities are also prominent: in London, 88\% of \ac{bss} members identified as white compared with 55\% of the general population~\cite{Fishman2016}, and in Washington, D.C., only 3\% of \textit{Capital Bikeshare} members were African-American and accounted for just 8\% of all cyclists in the area~\cite{Buck2013}, despite representing around half of the city’s population at the time~\cite{Fishman2016}. Across many contexts, subsequent studies have confirmed the persistence of these patterns~\cite{Shaheen2014, Fishman2013, Fishman2014, Bauman2017, Ogilvie2012, Steinbach2011, Willberg2021, Waldner2025}, even leading to some authors~\cite{Hoffmann2016} to describe bike sharing as “one of the most inequitable forms of sustainable transportation infrastructure.”

    \subsection{Gender patterns}

    Gender imbalances have often been attributed to differences in perceived risk and infrastructure design. Women tend to be more risk-averse than men~\cite{Garrard2008, Garrard2012}, and the lack of dedicated cycling infrastructure in many early \acp{bss} likely discouraged their participation. Nevertheless, some evidence suggests that \acp{bss} attract a higher proportion of women than regular cycling. Buck et al.~\cite{Buck2013} found that, compared with general cyclists in Washington, D.C., bike-sharing  users were more likely to be female, possibly because the bike-sharing environment reduces perceived risk—a result also observed in Sydney and Melbourne~\cite{Pucher2011}. Moreover, in London, women were more likely to make trips starting or ending in parks~\cite{Goodman2014}, suggesting a preference for routes away from motorized traffic and a tendency toward recreational rather than commuting trips, consistent with broader evidence that women favor traffic-free riding~\cite{Johnson2010}.

    Interestingly, gender imbalance in \acp{bss} is not universal, and cultural habits around cycling help explain much of the cross-national variation. In countries with traditionally low cycling rates such as the U.S., U.K., and Australia, men account for 65–90\% of total overall bicycle trips. By contrast, in places with long-established cycling cultures, such as Germany or Denmark, gender parity is common~\cite{Pucher2008}, and women even cycle more than men in Belgium~\cite{Witlox2004} and the Netherlands~\cite{Harms2014}. This cultural pattern extends to \acp{bss}: cities where cycling is male-dominated tend to see similar gender imbalances in bike-sharing usage~\cite{Fishman2016, Willberg2021}. Conversely, large dockless systems in China show little or no gender imbalance~\cite{Chen2020_exploring}, consistent with the country’s long history of widespread bicycle ownership and everyday cycling~\cite{Zhang2015}.

    \subsection{Socioeconomic and ethnic disparities}

    Socioeconomic disparities have largely been linked to the spatial distribution of the systems themselves and the requirements to use them. Because most \acp{bss} were initially deployed in central areas with high population density and economic activity, residents of peripheral or lower-income neighborhoods faced limited accessibility~\cite{Ogilvie2012, Goodman2014, Willberg2021}. Moreover, the requirement of a credit or debit card to register or rent bicycles has been seen to be an additional barrier for low-income individuals~\cite{Ricci2015, Ogilvie2012}. Moreover, pricing policies can reinforce these inequalities. In London, for example, Goodman and Cheshire~\cite{Goodman2014} found that a 2013 fare increase led to a decline in casual use among residents of the most deprived areas, illustrating how cost structures may further constrain participation among lower-income groups.

    However, these disparities are not inherent to bike-sharing. Emerging evidence shows that user profiles can shift when structural barriers are reduced. While many early evaluations in North America and Europe described \acp{bss} users as predominantly white, higher-income men, more recent findings challenge this assumption. A U.S. nationwide study revealed that African-American and Hispanic individuals were more likely to use bike-sharing than white users~\cite{Yu2024}, and that low-income individuals (earning under \$15,000) had a higher propensity to use \acp{bss} than higher-income groups. These results suggest that, under supportive conditions, bike-sharing can broaden mobility access for groups traditionally underserved by other transport modes—though this depends strongly on system design.

    \subsection{Age profile and inclusion of older adults}

    Another important question is whether the age profile of users has shifted over the years. The first wave of \ac{bss} users from the late 2000s are now middle-aged or older, but data indicate that bike-sharing still tends to be used far more by younger adults than by seniors. Usage declines with increasing age, and older individuals remain less likely to use bike-sharing than younger people~\cite{Yu2024}. In Helsinki, for instance, people older than 60 years were markedly underrepresented among bike-sharing riders, even after several years of operation~\cite{Willberg2021}. 

    Nevertheless, there are signs that older adults are participating in \acp{bss}. In Washington's D.C.'s \textit{Capital Bikeshare}, in 2018, about 24\% of its riders were aged 45 or over~\cite{Walljasper2018}, in Logroño's \textit{BiciLog}, men between 60 and 69 years old had almost the same levels of usage compared with 20 year-old men~\cite{CortezOrdonez2023}, and a similar pattern was found in València's system~\cite{PansSancho2023}. Moreover, in Madrid's \textit{BiciMAD} all group of users over ages 35 have been reported increase from 2015 to 2019~\cite{Julio2022}, and the authors suggest that it is due to the users' aging and the increasing acceptance of cycling in the society. Some of these may be early adopters who have continued using the service as they aged, while others might be newer members drawn from older age brackets. 

    In the last years, initiatives to engage older adults—such as outreach programs and educational group rides for seniors—have been reported in several cities~\cite{Ghisallo2023, BikeLNK2025}, and the growing availability of \acp{e-b} in bike-sharing fleets could encourage older ridership by reducing the physical effort of long or hilly trips. Recent evidence from the Netherlands shows that \acp{e-b} have led to higher usage among older women, with the gender gap narrowing in older generations. Among younger cohorts, however, the gap has widened, suggesting that age and gender interact in complex ways in \ac{e-b} adoption~\cite{Huang2024}.

    \subsection{Vehicle ownership and mobility context}

    Vehicle ownership plays a nuanced role in shaping the adoption of \acp{bss}, with context such as infrastructure and trip purpose significantly influencing usage patterns. Multiple studies have found an inverse relationship between private vehicle ownership and bike-sharing usage. In general, individuals without access to a car are more likely to use \acp{bss} than those with a car. Two U.S.-based studies provide clear evidence of this pattern. First, Mohiuddin et al.~\cite{Mohiuddin2024} reported that “frequent” bike-sharing users were often individuals without car access, who relied on a combination of biking and \ac{pt} for their daily mobility. Then, Yu~\cite{Yu2024} identified a statistically significant negative correlation between household vehicle ownership and cycling frequency, where households with more vehicles tended to make fewer bike trips. Similarly, Zhou et al.~\cite{Zhou2025} observed a negative and non-linear relationship between household vehicle ownership and bike-sharing usage in Shenzhen, China. Comparable findings have also been reported in Europe, where, for instance, bike-sharing users in Oslo were less likely to own a car than non-users~\cite{Mouratidis2022}.
    
    Further evidence from operational systems reinforces this trend. In Washington, D.C., more than half of \textit{Capital Bikeshare} members reported having no access to a private vehicle~\cite{CapitalBikeshare2013}. Notably, 5\% of users said they had sold a car after joining the system~\cite{CapitalBikeshare2013}, indicating that \acp{bss} can actively contribute to reducing car dependency. In fact, in a longitudinal study of Boston, Basu and Ferreira~\cite{Basu2021} found that the introduction of a new bike-sharing station led to a 2.2\% reduction in household vehicle ownership. Interestingly, a study in Montreal found that users holding a driver’s license had 1.5 times greater odds of using the system than those without it~\cite{BachandMarleau2012}. 

    At the same time, car ownership does not necessarily discourage bike-sharing use. Some evidence even shows car owners among the most active \ac{bss} users. A survey in Hangzhou, China found that \ac{bss} members had a higher rate of auto ownership (22\%) than non-members (11\%)~\cite{Shaheen2011_worldwide}. This may reflect the role of \acp{bss} as a time-saving alternative for short trips or a way to avoid traffic and parking constraints in congested areas. Likewise, in the Oslo study~\cite{Mouratidis2022}, although car access was initially associated with lower bike-sharing use, this effect disappeared after controlling for urban form and transit access. The findings point to a key nuance: it is not vehicle ownership alone that discourages bike-sharing use, but rather the spatial and infrastructural context in which that ownership exists. In compact cities with robust bike infrastructure and multimodal options, even car owners often find \ac{bss} convenient and efficient.

    Motorcycle ownership can play a similar role. In some regions, particularly in South Asia, two-wheel motorbikes dominate urban travel and compete directly with cycling. Yet, recent evidence from India reveals that \ac{bss} can attract even motorbike owners. A multi-city analysis found a notable mode shift from private two-wheelers to bike-sharing, a pattern described as “unique to India”~\cite{TUMI2023}. In these cases, \ac{bss} users cited availability and travel time savings as key factors for switching to shared bicycles, especially in heavily congested urban settings. Therefore, dense station networks and dockless bikes can encourage mode shift even away from motorized two-wheelers.

    Personal bicycle ownership, by contrast, does not always translate into bike-sharing use. Several studies indicate that individuals who already own a private bicycle may be less likely to use \ac{bss} regularly. For instance, Bachand-Marleau et al.~\cite{BachandMarleau2012} found that in Montreal, personal bike owners were underrepresented among regular \ac{bss} users. Similarly, in Washington, D.C., only 29\% of \textit{Capital Bikeshare} users reported owning a personal bicycle~\cite{Buck2013}.

    At the same time, other studies highlight that a lack of personal bike ownership is often what drives \ac{bss} usage in the first place. Factors such as concerns over theft, limited home storage, or the inconvenience of transporting a personal bicycle across different modes of travel make shared bikes more appealing~\cite{Teixeira2023, Ricci2015, Molina-Garcia2013, BachandMarleau2012}. In Lisbon, users described shared bicycles as "easier" and less stressful than owning their own~\cite{Teixeira2023}. Additionally, some users report that bike-sharing served as an entry point into cycling, as their initial experience with \ac{bss} encouraged them to later purchase a personal bicycle~\cite{Fishman2016}.

    Together, these findings highlight that while private vehicle ownership—whether car, motorcycle, or bicycle—can shape the likelihood and manner of bike-sharing use, it is often the broader context of infrastructure, convenience, and perceived barriers that determines actual user behavior.

    \subsection{Psychographic attributes}

    Beyond demographic characteristics, a growing body of research highlights the role of psychographic attributes—that is, the attitudes, motivations, perceived benefits, and social norms that shape how individuals evaluate and engage with \acp{bss}. These factors help explain why some groups adopt bike-sharing while others do not, even when physical access and pricing are similar.

    At the individual level, studies using behavioural frameworks show that positive attitudes toward cycling, perceived benefits (including health, exercise and environmental advantages), and a sense of confidence and control in traffic (perceived behavioural control) are significantly associated with intentions to use bike-sharing among college students~\cite{Chen2022_Predicting}. Beyond these individual attitudes, research comparing users and non-users finds that mobility-related identities and lifestyle orientations also matter: individuals who hold more mobility-oriented and bicycle-friendly lifestyles are more likely to become users, whereas those who value car-based mobility or place importance on traditional transport preferences tend to avoid \acp{bss}~\cite{Stockle2020}. Recent segmentation analyses further illustrate how these lifestyle orientations translate into distinct usage patterns: some users appear in a multimodal, efficiency-oriented segment, others in a multi-purpose usage segment that includes commuting and social trips, while more occasional users—characterized by higher income and car ownership—engage with bike-sharing at lower frequency~\cite{Mohiuddin2024}.

    At a broader cultural level, the meaning of cycling within a city’s mobility culture shapes the conditions under which bike-sharing is perceived as a legitimate and desirable mobility practice. Rather than being determined solely by infrastructure or pricing, the position of cycling within local mobility cultures reflects dominant values, norms, and ideas about what mobility should represent~\cite{Nikolaeva2019}. Where mobility is framed around flexibility, shared access and collective benefit, cycling and bike-sharing are more readily integrated into everyday mobility; where mobility remains strongly associated with speed, autonomy and automobility, alternative modes may struggle to gain symbolic and practical legitimacy. Complementing this, survey-based studies show that bike-sharing participation varies across value-based lifestyle milieus: groups oriented toward flexibility, innovation, and urban mobility tend to use bike-sharing more, while groups with more traditional or precarious orientations show lower uptake~\cite{Borgstedt2019}. Relatedly, peer acceptance and perceived social legitimacy strongly influence participation, with higher usage where bike-sharing is seen as a familiar and supported practice~\cite{Boenigk2019, Georgi2019}.


\section{Determinants of system usage}

This Section examines the determinants of \ac{bss} usage across multiple scales. It begins by outlining the operational indicators used to observe and evaluate system performance in practice. It then discusses how usage is shaped by broader city-level contextual conditions, system-level design choices, and neighbourhood- and station-area built environment characteristics.

    \subsection{Operational indicators}
    \label{sec:operational_indicators}

    Operators of \acp{bss} typically rely on a set of operational indicators to monitor and sustain system performance in real time. Such monitoring is essential because the success of a bike-sharing scheme is not guaranteed; several high-profile implementation failures~\cite{Yle2015_HelsinkiBikeShare} have shown how operational problems can quickly undermine public confidence. In response, many operators now use well-defined performance indicators to support continuous monitoring, proactive management, and evidence-based adjustments.
    
    System-wide utilization is commonly assessed through two complementary indicators: the average number of \ac{tdb} (or turnover rate) and the average daily trips per 1,000 residents~\cite{Waldner2025, Yanocha2018, Yahya2017}. The \ac{tdb} rate measures how intensively the fleet is being used, calculated as the total number of daily trips divided by the number of bicycles in the system. A broadly recommended benchmark ranges between 4 and 8 daily uses per bicycle~\cite{Yanocha2018}, as rates below this threshold may jeopardize financial sustainability, while much higher rates can lead to shortages during peak demand. For example, \ac{tdb} of 6.4 were reported for both Barcelona and New York, and 5.9 for Guangzhou in 2017~\cite{Yanocha2018}, while a cross-city comparison of 75 systems found values ranging from 0.22 to 8.4~\cite{deChardon2017}.

    The average daily trips per 1,000 residents is a penetration indicator that reflects the degree to which the system is integrated into everyday mobility patterns within its service area. Most cities report values below 30~\cite{Yanocha2018}, although cases such as Mexico City (\textasciitilde100) and Dublin (\textasciitilde75) exhibit substantially higher figures due to large numbers of commuters who use the system but do not reside within the service area.

    Together, these utilization indicators offer a broad picture of overall demand and fleet efficiency. However, their interpretation depends on more structural characteristics of the system~\cite{Yanocha2018}. These include the size of the service area (typically defined as the union of 500-meter service radii around stations), the population within that service area, the total fleet size, the number of stations, and the total number of docks. Moreover, more sophisticated indicators such as the percentage of low-income populations who live and/or work within the service area are also used to assess equity goals~\cite{Yanocha2018}.

    Beyond system scale, service reliability is assessed using indicators that capture the operational state of stations and bicycles. A key measure is the percentage of time stations are non-operational, meaning they are either completely empty (no bicycles available) or completely full (no parking spaces available), as both situations prevent intended use and directly affect users’ ability to complete their trips~\cite{Santander2022, Rossow2018}. Relatedly, operators often examine average dock occupancy to identify spatial or temporal imbalances and to inform rebalancing efforts.
    
    Maintenance performance is also closely monitored~\cite{Rossow2018, Santander2022} as it is crucial for the long-term sustainability of the system. Typical indicators include the bike's lifespan, the share of the fleet temporarily out of service or the breakdown rate of the bike's components. Moreover, the rise of \acp{e-b} introduces the necessity to add a new kind of monitoring: the state of charge of batteries across the fleet.

    Finally, operators track financial efficiency to assess long-term sustainability~\cite{Yanocha2018}. Common indicators include capital and operating cost per bike, which support budgeting and benchmarking, and operating cost per trip, which reflects how efficiently the system delivers service as it scales.

    \subsection[Factors influencing BSS usage]{Factors influencing \acs{bss} usage}
    \label{sec:factors_influencing_bss_usage}

    Understanding why people use \acp{bss} requires considering multiple layers of influence, ranging from system design decisions to the characteristics of the built environment and broader city-wide conditions. According to OBIS framework~\cite{OBIS2011}, these determinants can be grouped into endogenous factors related to the system’s configuration and operation, and exogenous factors linked to the urban and environmental context. In this Section, the key determinants of \ac{bss} usage are discussed at the city, system, and station scales.


        \subsubsection{Contextual factors}

        Contextual factors are predominantly exogenous, meaning they are not directly controlled by the system operator. One basic contextual factor is the size of the city in which the \ac{bss} operates. City size is typically measured by resident population, although the presence of non-resident users (such as work travellers or tourists) also can contribute to the potential demand base when the BSS is open to non-residents~\cite{Ferrando2007}. In general, larger cities provide a broader market for bike sharing, and several guidelines recommend a minimum population of around 200{,}000 inhabitants for a system to be viable~\cite{Ferrando2007, CTCN2017}. However, successful implementations in smaller cities demonstrate that this is not a strict threshold. For example, Boulder, U.S. has maintained a functioning \ac{bss} with a population of approximately 108,000 residents~\cite{Yanocha2018}. Beyond overall population, both population and employment densities are consistently associated with higher usage, as they concentrate trip origins and destinations~\cite{Tran2015, Chen2021_Identifying}.

        Cycling culture and modal split also play an important role. In cities where cycling is already a familiar and socially accepted mode of transport, bike-sharing tends to be adopted more readily and used more frequently~\cite{Fishman2016, Willberg2021}. Investments in safe cycling infrastructure further amplify this effect. For example, in New York City, the installation of protected bike lanes near bike-sharing stations increased monthly trips at those stations by approximately 18\%, while simple painted lanes generated a surprisingly 14\% increase~\cite{Moran2025}. Thus, \acp{bss} tend to perform better where the street environment is perceived as safe and inviting.

        % GUON2022: First, cities in West Europe, East Asia, and North America \& Oceania are generally cycling-oriented, transit-oriented (or transit-dependent), and automobile-oriented, respectively (26, 31-33).
        % Ampliar amb secció GUON2022:Variance of The Built Environment Effects Across Mobility Culture

        Topography is another critical determinant. Steep slopes discourage cycling: positive gradients above 4\% tend to reduce comfort, and those above 8\% are often actively avoided~\cite{Midgley2011}. Accordingly, several cycling design guidelines recommend avoiding infrastructure on slopes greater than 10\%~\cite{AASHTO1999bicycle, CycleHighways2025, CROW2006}. Cities with pronounced topography may therefore experience spatially uneven demand unless mitigated through \acp{e-b}, and adapted rebalancing operations may be needed to ensure equitable access.

        Finally, weather strongly shapes daily and seasonal ridership. A multicity analysis of nearly 100 million trips across 40 cities found strong dependencies on temperature and precipitation~\cite{Bean2021}. Usage increases with warmer temperatures up to approximately 27–28 °C, beyond which it declines~\cite{Julio2022}. Rainfall and snow cause sharp decreases~\cite{Bean2021, Wang2021, Julio2022}, particularly among women~\cite{Nahal2018}, and winter conditions can lead to deep seasonal troughs in temperate climates~\cite{Julio2022}. Broader evidence also shows that favourable weather—such as dry conditions, mild temperatures, and hours of sunshine—positively affects bike-sharing use~\cite{Corcoran2014, FaghihImani2014, Gebhart2014, Caulfield2017, Campbell2016}, whereas high temperatures, humidity, wind speed, and poor air quality have negative effects~\cite{Campbell2016, Corcoran2014, Caulfield2017}. Nonetheless, \acp{bss} operate successfully in diverse climates~\cite{OBIS2011, Yanocha2018}, with some cities adapting operations, such as winter closures in Oslo, Stockholm, and Montreal~\cite{OBIS2011}.

        
        \subsubsection{System-level and operational factors}

        The overall configuration of a \ac{bss} has a strong influence on its usage. A fundamental determinant is the extent and shape of the service area. Systems that cover dense, mixed-use neighbourhoods with high trip-generation potential--typically city centres--tend to exhibit higher ridership~\cite{Yanocha2018}. However, extending coverage to lower-density areas can also support demand when these areas lack adequate transport connectivity. Crucially, the service area must be large enough to encompass a meaningful share of residents’ origins and destinations; otherwise, the convenience of the system is reduced and usage is limited.

        Within the service area, the density and distribution of stations strongly condition how easily users can incorporate bike-sharing into their daily trips. Planning guidelines consistently identify these factors as key drivers of ridership~\cite{Yanocha2018, Ferrando2007, OBIS2011}. Typical station density values range between 10 and 16 stations per km$^{2}$~\cite{Yanocha2018, Ferrando2007, OBIS2011}, a level that generally ensures that most users are within 300–500 metres of a station. This is supported by Kabra et al.~\cite{Kabra2018}, who found that nearly 80\% of trips are originated within 300 metres of a station. Notably, users also tend to perceive systems with dense station networks and broad coverage as more reliable and successful~\cite{OBIS2011}.

        Adequate bicycle supply is equally important. Recommended fleet sizes commonly range between 10 and 30 bicycles per 1,000 residents, which supports target usage levels of approximately 4–8 \ac{tdb}~\cite{Yanocha2018}. Medium and large systems often provide a higher number of bikes and docks per station to ease rebalancing and maintain availability~\cite{OBIS2011}. Finally, the ratio of docks to bicycles influences parking convenience: a dock-to-bike ratio of roughly 2–2.5 helps ensure that users can find a free dock during peak arrival periods, reducing the likelihood of stations becoming full~\cite{Yanocha2018}.

        Another design choice is vehicle type. This is a key decision, as the incorporation of \acp{e-b} can fundamentally reshape how users experience physical effort. Evidence from Madrid’s fully electric BiciMAD system illustrates this clearly: user surveys show that the perceived importance of street slope has steadily diminished over time, indicating that electrification effectively neutralizes topographic barriers and allows riders to use the system with substantially less effort~\cite{Julio2022}. For this reason, \acp{e-b} have been shown to boost usage by lowering physical exertion and attracting new users~\cite{Fyhri2015}. This effect is particularly pronounced in hilly cities. In Park City, U.S., a mountainous city characterized by steep terrain, adding \acp{e-b} significantly increased ridership during weekends and summer months, when leisure-oriented trips are more frequent~\cite{He2019}.

        Finally, pricing also plays a meaningful role in shaping usage. Cost influences how often and by whom the system is used, as higher per-trip or access fees can discourage casual and occasional users. Evidence from London shows that the 2013 price increase in its \ac{bss} led to a reduction in casual trips, particularly among residents in more deprived areas~\cite{Goodman2014}.


        \subsubsection{Station-level spatial and built environment factors}

        Beyond context and system design, two stations of the same \ac{bss} can perform very differently due to the immediate local context. At the station scale, usage reflects the interaction between multiple factors: those concerning the built environment surrounding the station and those concerning the station design~\cite{Guo2022, Guo2020, Liu2019, Ma2020_comparison, Chen2021_Identifying}. These factors shape not only the magnitude of trips but also their temporal profile and the balance between arrivals and departures.

        The built environment refers to the physical setting that frames human activity~\cite{Guo2022}. It typically comprises three interrelated components~\cite{Handy2002}: (i) urban design, which concerns the arrangement and appearance of physical elements such as street layout, block shapes, and tree cover; (ii) land use, which describes how activities are distributed across space, including residential, commercial, office, industrial, and other functions; and (iii) the transportation system, which includes roads, sidewalks, bike lanes, and \ac{pt} facilities, as well as their level of service, such as traffic conditions and \ac{pt} frequency.

        The relationship between the built environment and \ac{bss} usage has received growing attention over the past decade, though findings vary across cities and mobility cultures~\cite{Guo2022}. Table~\ref{tab:built-environment-factors} summarizes commonly studied station-level built environment factors.

        \begin{table}[hpbt!]
            \caption[Commonly studied built environment factors influencing BSS usage]{
                \textbf{Commonly studied built environment factors influencing \ac{bss} usage.} 
                Adapted from Guo et al.~\cite{Guo2022}.
            }
            \label{tab:built-environment-factors}
            \centering
            {\footnotesize 
            \begin{tabular}{lll}
                \toprule
                Category & Sub-category & Measurements \\
                \midrule
                Land use & Land use type & Share or area of residential land \\
                & & Share or area of office land \\
                & & Share or area of industrial land \\
                & & Share or area of commercial land \\
                & & Share or area of green/open space \\
                & Land use mix & Land use entropy index \\
                & 	\acp{poi} & Number of shopping malls \\
                & & Number of universities or schools \\
                & & Number of parks \\
                & & Number of recreation facilities \\
                & & Number of restaurants \\
                & & Number of retail stores \\
                \midrule
                Transportation system & Road network & Length of bicycle lanes \\
                & & Length of major roads \\
                & & Length of minor roads \\
                & & Length of highways or regional roads \\
                & & Length or presence of paved trails \\
                & & Number of street intersections \\
                & \ac{pt} infrastructure & Number of subway or rail stations \\
                & & Length of subway or rail lines \\
                & & Number of bus stops \\
                & & Number of bus lines \\
                \midrule
                & Bike-sharing infrastructure & Number of stations \\
                & & Station dock capacity \\
                \midrule
                Urban design & Amenities & Presence of street trees or shading \\
                & & Presence of street lighting \\
                & Accessibility & Distance to city hall or government centers \\
                & & Distance to central business district \\
                & & Distance to nearest \ac{pt} stop \\
                \midrule
                Urban form & Density & Population or household density \\
                & & Employment or job density \\
                \bottomrule
            \end{tabular}
            }
        \end{table}

        Land-use composition around a station consistently emerges as one of the main predictors of use. Across many cities, stations in dense residential and commercial areas tend to register higher trip volumes, while those in purely office or green-space contexts exhibit fewer departures and returns~\cite{FaghihImani2014, FaghihImani2016, MateoBabiano2016, FaghihImani2017, Sun2018, FaghihImani2017Empirical}. In Montreal, for instance, stations near the city center and surrounded by restaurants, jobs, and mixed land uses showed markedly higher ridership~\cite{FaghihImani2014}. Likewise, in Minneapolis–St. Paul, job clusters and food-related businesses strongly correlated with station activity~\cite{Wang2016}. Yet multi-city analyses indicate diminishing returns once local densities exceed approximately 20,000 residents/km$^{2}$ and 12,000 jobs/km$^{2}$~\cite{Chen2021_Identifying}. Density effects also vary geographically: in Western Europe and North America, demand typically peaks in dense urban cores, whereas in East Asian megacities such as Seoul, Beijing, Taipei, and Shanghai, excessive density and pedestrian congestion can deter cycling by raising perceived risk and discomfort~\cite{Liu2019, Ji2020, Rixey2013, Lin2020}.

        Beyond functional density, specific land uses attract distinctive trip purposes. Parks, waterfronts, and leisure \acp{poi} (e.g., cinemas, heritage sites, tourist landmarks) often generate recreational or loop trips that start and end at the same station, especially within large parks~\cite{Tran2015, MateoBabiano2016}. Moreover, being near to one of these amenities is not all, the closer a station is, the greater its expected usage~\cite{Wang2016}.

        Proximity to \ac{pt} is another key determinant. \acp{bss} interact closely with metro, rail, and bus networks, as nearby transit stations typically increase both departures and arrivals~\cite{Tran2015, Lee2021}. In China, dockless systems show particularly strong feeder relationships with metro and bus services, thanks to their flexible parking and broader reach into residential and industrial areas~\cite{Guo2020, Guo2021Dockless, Ji2020}. Some authors suggest that optimal distances for bike–metro transfers range between 0.8 and 1.5~km~\cite{Zhao2017}, though the strength of this relationship varies by city morphology and transit design~\cite{Lin2018}. Overall, the distribution of metro and rail stations and their catchment areas remain a fundamental determinant of first- and last-mile demand.

        Street context further conditions station use. The presence of bike lanes increases willingness to ride and correlates positively with both arrivals and departures~\cite{ElAssi2017}. Stations located along minor streets or near intersections typically show higher activity, while those adjacent to highways, railways, or major arterials register lower usage~\cite{FaghihImani2014, FaghihImani2016, FaghihImani2016SpatioTemporal, ElAssi2017}. Moreover, local amenities that improve comfort and safety such as street trees, shading, and lighting also exhibit positive associations with departures~\cite{Liu2019}.

        Topography introduces another spatial asymmetry. Stations at higher elevations generally record fewer arrivals and more departures, as riders tend to avoid uphill trips and favor downhill returns~\cite{Morency2015, MateoBabiano2016}. However, the growing prevalence of \acp{e-b} can mitigate this effect.

        Finally, the design and capacity of the stations themselves influence ridership. Dock capacity—and, by extension, available bike supply and the dock–bike ratio—are consistently identified as strong predictors of both departures and arrivals~\cite{ElAssi2017, Yang2020, Zhao2021}. Moreover, station visibility and accessibility are highly recommended for attracting trips~\cite{Ferrando2007, Yanocha2018, CROW2006}.

        
\section[Riding behaviour in BSSs]{Riding behaviour in \acp{bss}}

    \subsection{Trip characteristics and user behavior}


% ############################# Purpose #############################
Although individual \acp{bss} exhibit diverse spatial and temporal patterns, empirical research consistently reveals marked daily and weekly regularities in usage~\cite{Fishman2016}. Most trips are utilitarian, including commuting and errands, but leisure-oriented trips such as tourism and recreational rides are also common~\cite{Ricci2015, Fishman2016, Willberg2021, CapitalBikeshare2013, Xu2019_Unravel, Buck2013, Raux2017, Lin2018}. In Lyon, for instance, survey data from the \textit{Velo’v} system show that commuting is the main purpose, accounting for 36\% of weekday trips among annual members~\cite{Raux2017}. Leisure and shopping trips represented 18\% and 12\%, respectively, confirming that while work-related mobility dominates, a substantial share of rides is also motivated by non-commuting activities.

% ############################# Trips characteristics #############################

While trip characteristics can vary depending on the specific city in which the \ac{bss} operates, existing research consistently reports similar magnitudes across contexts. On average, trips typically cover between 1 and 2 kilometers, last 10–20 minutes, and occur at mean speeds of roughly 8–15~km/h~\cite{Fishman2014, Zhang2016Bicycle, Levy2019, Zhao2015, MateoBabiano2016, FaghihImani2017_Hail, Yang2024}. In most Western systems, about 95\% of users return the bike within 30 minutes~\cite{Zhao2015, MateoBabiano2016}, aligning with the common free-ride threshold, and trip lengths are typically shorter than those made on private bicycles~\cite{CastilloManzano2015}. 

Trip duration and distance are shaped by temporal, spatial, and user factors. Travellers usually take shorter, slower rides on weekends than on weekdays, particularly for leisure-oriented purposes~\cite{MateoBabiano2016}. Casual or day-pass users tend to ride longer than annual members~\cite{Buck2013}, reflecting their more recreational motivations. However, because commuting dominates \ac{bss} activity, such differences remain modest at the aggregate level. In Helsinki, for instance, 82\% of users were annual subscribers who accounted for 92\% of total trips, confirming that regular members generate most system usage~\cite{Willberg2021}.

Usage frequency among registered members nonetheless varies widely. In London, nearly half of all members reported not having used the service during the previous month~\cite{Fishman2016}. Similar proportions were found in Washington, D.C., where 21\% of female and 13\% of male members reported no rides in a typical month~\cite{Buck2013}, and in Australia, where 46\% of annual members recorded no trips at all in the prior month~\cite{Fishman2014}. Comparable heterogeneity has been observed elsewhere: in Lyon and Vancouver, most users were irregular riders, while a small share of ‘super-users’—often men and young adults—accounted for a disproportionate number of trips~\cite{Vogel2014, Winters2019}. These findings suggest that many subscribers view bike sharing as an occasional complement to other transport modes rather than a primary means of travel, often maintaining memberships to retain flexibility or support the system~\cite{Fishman2016}.

% ############################# e-bikes and dockless systems #############################

The introduction of \acp{e-b} has further reshaped these mobility patterns. Compared with \acp{m-b}, \ac{e-b} trips tend to be longer, faster, and spatially more dispersed~\cite{Yang2024}. Electrification therefore extends the effective catchment area of bike-sharing systems, enhances accessibility in hilly zones, and broadens the demographic base of potential users. Evidence from Madrid’s BiciMAD confirms that street slope has become a comparatively minor constraint, while flexibility and environmental awareness have emerged as dominant motivators for use~\cite{Julio2022}. Moreover, comparative analyses show that docked systems display clearer commuting peaks, whereas dockless services exhibit more temporally distributed and recreation-oriented profiles~\cite{Bonilla-Alicea2020}.

    \subsection{System-level spatio-temporal patterns}

    Temporal regularities are among the most consistent features of \acp{bss}. On weekdays, ridership peaks during the morning (6–9 a.m.) and evening (4–7 p.m.) rush hours, reflecting work-related mobility~\cite{Fishman2016, Zhao2015}, whereas on weekends demand shifts toward the middle of the day~\cite{Fishman2016, CortezOrdonez2024}. Because commuting dominates weekday activity, \ac{tdb} values are also higher on working days~\cite{FaghihImani2014, ElAssi2017}. Seasonal patterns further modulate demand: in cities with cold winters such as Helsinki or Chicago, usage rises during warmer months and declines sharply in winter~\cite{Willberg2021, Zhou2015, Wang2021}, while in milder climates like Barcelona activity is reduced during the hottest summer months~\cite{CortezOrdonez2024}. Moreover, precipitation substantially reduces trip volumes~\cite{Wang2021}.

    Beyond temporal rhythms, \acp{bss} display distinct spatial structures that mirror the organization of urban space. Across cities, activity concentrates in dense, mixed-use, and transit-accessible areas, while peripheral zones show lower and more irregular demand~\cite{FaghihImani2017, Willberg2021}. Trips commonly originate and terminate near residential, employment, or major public-transport hubs, reflecting both the spatial configuration of the station network and surrounding land-use composition. In Helsinki, for example, 79\% of all bike-share trips were made by residents living within the 40\% population-coverage area of the system, underscoring that proximity to stations strongly shapes demand~\cite{Willberg2021}. This finding aligns with results from Montreal, where individuals residing within 500 m of a docking station were more than three times as likely to use the service~\cite{BachandMarleau2012}. Flows often align with commuting corridors~\cite{Willberg2021}: users living beyond the main service area frequently begin trips at major transit hubs in the morning and return in the afternoon, integrating cycling with public transport.

    \ac{od} analyses across several cities reveal that most rides occur within single functional zones—typically home–work relations or neighborhood loops. In Richmond, U.S., for example, the majority of trips remained confined within either business or residential areas, highlighting the localized and land-use-dependent nature of \ac{bss} demand~\cite{Yang2024}. Travel-pattern differences also relate to users’ home areas: in Minneapolis–St Paul, residents of minority and lower-income neighborhoods used the BSS more frequently and across a more diverse range of times and places compared with other users~\cite{Wang2019}.

    These global regularities are mirrored in specific urban contexts. In Taipei, Liu and Lin~\cite{Liu2019} identified four recurrent spatial patterns of station activity using cluster analysis. Stations in the city center and along metro corridors exhibited high dock turnover, multiple daily peaks, and longer \ac{od} distances, characteristic of mixed or commercial districts dominated by work-related travel. In contrast, peripheral stations showed lower turnover and longer trips, reflecting occasional use from residential areas. On weekends and holidays, demand shifted toward stations near leisure amenities and metro-accessible destinations, indicating a transition to recreation-oriented mobility. Likewise, in Tel Aviv, bike-sharing activity clustered strongly in the dense urban core, forming a pronounced “usage ring” around the city center and along the waterfront~\cite{Levy2019}. Short trips concentrate in central areas and often complement public transport, whereas longer rides follow north–south and northeast–southwest axes aligned with major cycling corridors, suggesting both substitution effects and the spatial influence of dedicated infrastructure.

    % ############################# Spatial patterns nuances #############################

    Finally, temporal profiles also reflect the functional purposes embedded in nearby land uses. Stations near universities typically show strong morning arrival peaks but weak departure peaks, as users converge at similar times for classes but leave at varying hours~\cite{FaghihImani2014}. In contrast, stations adjacent to restaurants and entertainment venues show negligible morning activity but pronounced evening peaks aligned with opening hours~\cite{FaghihImani2016}. Employment and job density likewise elevate overall station activity, mirroring the temporal structure of commuting flows~\cite{Tran2015, Chen2021_Identifying}.
